\chapter{Per-Stock Model B Results}
\label{AppendixB}

Table \ref{tab:perstockB} shows the per-stock $R^2$ values for Model B (without autoregressive input), highlighting the variation in feature-driven predictability across sectors.

\begin{table}[h!]
\caption{Per-Stock $R^2$ for Model B (Feature-Only)}
\label{tab:perstockB}
\centering
\begin{tabular}{llc}
\hline
\textbf{Stock} & \textbf{Sector} & \textbf{$R^2$ (Model B)} \\
\hline
ADANIPORTS & Infrastructure & 0.503 \\
HAL & Defence & 0.441 \\
COALINDIA & Energy & 0.387 \\
TITAN & Consumer & 0.352 \\
TATASTEEL & Metals & 0.298 \\
LT & Infrastructure & 0.276 \\
RELIANCE & Energy & 0.245 \\
WIPRO & IT & 0.221 \\
TCS & IT & 0.198 \\
INFY & IT & 0.156 \\
HDFCBANK & Banking & 0.023 \\
AXISBANK & Banking & $-$0.412 \\
ICICIBANK & Banking & $-$0.876 \\
SBIN & Banking & $-$1.525 \\
\hline
\end{tabular}
\end{table}

The results reveal a clear sectoral pattern: banking stocks are the least predictable from exogenous features alone, while infrastructure and defence stocks show the highest feature-driven predictability. This pattern may reflect the greater sensitivity of banking stocks to monetary policy signals and RBI announcements that are not captured in the current feature set.
